\section{Introduction}
\label{sec:introduction}

Personal identity documents --- passports, driving licences,
national identity cards --- are among the most sensitive
artefacts in modern administrative life.
They are routinely scanned, emailed, and stored by employers,
landlords, banks, and government agencies, creating large
corpora of high-value personally identifiable information
(PII) that are difficult to redact reliably and even
harder to audit credibly.
Existing redaction tools fall into two categories, each
with a fundamental limitation: cloud-based APIs (AWS
Comprehend, Google Cloud DLP) process text only and
transmit sensitive documents to third-party servers,
violating the very privacy they claim to protect;
and desktop tools (PDF-XChange, Adobe Acrobat) require
manual identification of sensitive regions, are error-prone
under time pressure, and produce no verifiable audit trail.

The inadequacy of current approaches is not merely
theoretical.
A nurse who must redact a scanned medical ID before
forwarding it has no tool that is simultaneously
\emph{accurate}, \emph{local}, and \emph{auditable}.
An organisation subject to GDPR Article~25 (data protection
by design) cannot demonstrate compliance with a manual
highlighter.
A legal firm handling identity documents in bulk cannot
afford the latency and privacy exposure of a cloud
round-trip for every page.
These gaps motivate a different design point: a system that
runs entirely on edge hardware, understands documents
visually rather than as raw text, and produces a
cryptographically verifiable record of every redaction
decision.

We present \textbf{Omni-Shield}, the first local-first
multi-modal PII redaction system that combines a
fine-tuned vision-language model (VLM) with a nine-agent
orchestration pipeline, a blockchain audit trail,
and real zero-knowledge proofs --- running entirely on
a consumer-grade GPU (NVIDIA RTX~5060, 8\,GB VRAM).
No document leaves the device.
No cloud account is required.
Every redaction is anchored on-chain and accompanied by
a Groth16/BN128 ZK-SNARK that proves redaction occurred
without revealing what was redacted.

\subsection*{The core challenge}

Identity document redaction is harder than general-purpose
NLP-based PII detection for three reasons.

\textbf{Visual PII cannot be detected from text alone.}
Signatures and face photographs --- the two most legally
significant PII elements on identity documents --- are
not text.
They cannot be found by regex, spaCy NER, or any
text-only pipeline regardless of model size.
Our evaluation confirms this: Microsoft Presidio, a
mature NLP-based system, achieves F1\,=\,0.000 on
signatures and faces across 150 test documents.

\textbf{Layout determines meaning.}
On a passport, the string \texttt{15\,MAR\,1987} is a
date of birth.
On an invoice, the same string is an issue date and
should not be redacted.
A system that does not understand document layout ---
which field labels surround a value, where the value
appears spatially on the card --- will systematically
misclassify or miss PII.

\textbf{Redaction without verification is insufficient
for compliance.}
A redacted PDF is only trustworthy if a verifier can
confirm (a)~which document was redacted, (b)~when,
(c)~by whom, and (d)~that the redaction was performed
correctly.
Existing tools produce none of these guarantees.
A blockchain-anchored audit trail with ZK-SNARK proof
addresses all four.

\subsection*{Contributions}

This paper makes four contributions:

\begin{enumerate}

  \item \textbf{Multi-modal nine-agent redaction pipeline.}
  We design and implement a nine-agent pipeline
  (RouterAgent, DeskewAgent, OCRAgent, LayoutAgent,
  TextPIIAgent, ContextAgent, BBRefinerAgent,
  CriticAgent, RedactionAgent) that processes six
  document formats (images, PDFs, audio recordings,
  plain text, Word documents, PowerPoint presentations)
  on a single 8\,GB GPU at 4-bit NF4 quantisation.
  The pipeline achieves a mean end-to-end latency of
  10.42\,s for image redaction and under 50\,ms for
  text and PDF documents.

  \item \textbf{Knowledge distillation for document PII.}
  We fine-tune Qwen2-VL-2B-Instruct on 2,000 synthetic
  identity documents across six countries using QLoRA,
  with teacher labels generated by the Claude API.
  The fine-tuning teaches the model to emit consistently
  structured \texttt{LABEL\,|\,VALUE\,|\,BBOX} output
  required by the downstream pipeline --- a contribution
  to output format adherence rather than raw detection
  capability, which we measure and discuss honestly.
  As a notable emergent capability, face detection
  generalises from the training distribution without
  explicit face annotation.

  \item \textbf{Cryptographic audit trail with ZK-SNARK.}
  Every redaction produces a blockchain-anchored record
  (AuditLog.sol, Groth16/BN128 via ZoKrates) that proves
  redaction occurred without revealing the redacted values.
  Blockchain verification completes in 1.2\,ms via O(1)
  Solidity mapping lookup.
  ZK proof generation takes 20.3\,s but executes
  asynchronously, adding no user-facing latency.
  Each anchor costs 264,333 gas units --- approximately
  \$0.006 on Polygon L2.
  We also identify and fix a hash linkability
  vulnerability (SHA-256 without salt produces
  identical on-chain hashes for repeated documents)
  and implement a timestamp-salted commitment scheme.

  \item \textbf{Comprehensive evaluation on synthetic
  and real documents.}
  We evaluate Omni-Shield on 150 annotated synthetic
  identity cards and 19 real specimen documents
  (sourced from Wikimedia Commons).
  We compare against three baselines: Microsoft Presidio
  (rule-based NER, Macro F1\,=\,0.057), Llama~3~8B
  (LLM prompting, Macro F1\,=\,0.270), and our system
  (Macro F1\,=\,0.187 synthetic, 0.269 real).
  Omni-Shield is the only system to detect signatures
  (F1\,=\,0.870) and is the only baseline-outperforming
  system on real documents.
  We further report energy consumption
  (0.289\,kWh per 1,000 image redactions), carbon
  footprint under different grid assumptions, and
  the quantified ``privacy premium'' of local-first
  processing.

\end{enumerate}

\subsection*{Key findings}

Three findings stand out from our evaluation.
First, \emph{visual PII drives overall performance}:
ablation shows that disabling signature and face detection
reduces Macro F1 by 54\%, more than any other component.
Second, \emph{Omni-Shield generalises to real documents}:
Macro F1 on 19 real specimen cards (0.269) exceeds
performance on synthetic test cards (0.187), suggesting
the training distribution transfers well.
Third, \emph{text-only LLMs are complementary, not
competing}: Llama~3~8B outperforms Omni-Shield on
text-based PII (DOB F1\,=\,0.854, ID numbers 0.778)
but scores zero on signatures and faces --- confirming
that the multi-modal approach is necessary for complete
identity document redaction.

\subsection*{Paper structure}

Section~\ref{sec:related} surveys related work.
Section~\ref{sec:architecture} describes the system
architecture and nine-agent pipeline.
Section~\ref{sec:knowledge_distillation} covers the
synthetic dataset, knowledge distillation methodology,
and QLoRA fine-tuning.
Section~\ref{sec:evaluation} presents quantitative
evaluation against baselines and ablation studies.
Section~\ref{sec:crypto} details the cryptographic
audit trail, ZK-SNARK circuit, and blockchain design.
Section~\ref{sec:limitations} discusses limitations
and future work.
Section~\ref{sec:conclusion} concludes.
Code, evaluation harness, and the fine-tuned adapter
are available at \url{https://github.com/[anonymous]}.
